% ------------------------------------------------------------
% Chapter 4: Friedmann-Like Equations from Lamphrodynamics
% ------------------------------------------------------------

\chapter{Friedmann-Like Equations from Lamphrodynamics}
\label{chap:friedmann-like}

\section{Overview}

Classical cosmology relies on the Friedmann–Lemaître–Robertson–Walker
(FLRW) class of metrics, representing homogeneous and isotropic universes
with scale factor $a(t)$. The standard Friedmann equations follow from
Einstein’s field equations applied to the FLRW ansatz. In \rsvpabbr{}, the
modified Einstein equations
\[
  G_{\mu\nu} + \Lambda g_{\mu\nu} + \Theta_{\mu\nu} = 8\pi G\,T_{\mu\nu}
\]
lead, under homogeneity and isotropy assumptions, to \emph{modified}
Friedmann‐like equations with explicit entropic contributions. This chapter
derives these equations in detail and discusses physical interpretation,
including the effective equation of state and critical densities.

\section{Lamphrodynamic FLRW Ansatz}
\label{sec:flrw-ansatz}

We adopt the standard FLRW metric
\[
  ds^2 = -c^2dt^2 + a(t)^2\left(
    \frac{dr^2}{1 - kr^2} + r^2d\Omega^2
  \right),
\]
with curvature parameter $k\in\{-1,0,1\}$. Homogeneity and isotropy impose
that the lamphrodynamic fields $(\PhiField,\vField,\SField)$ respect FLRW
symmetry. We assume
\[
  \vField(t,\mathbf{x}) = 0,
\qquad
  \SField(t,\mathbf{x}) = \SField(t),
\qquad
  \PhiField(t,\mathbf{x}) = \PhiField(t),
\]
so that only time dependence remains. (More general anisotropic
configurations are possible but beyond the scope of this chapter.)

\begin{remark}[Physical interpretation]
In the FLRW context, lamphrodynamic entropy becomes a \emph{cosmic entropy
density} varying only with cosmic time. Entropy production and lamphrodynamic
couplings generate effective pressure and energy‐density terms modifying the
cosmological expansion.
\end{remark}

\section{Modified Einstein Tensor}

In the FLRW metric, the Einstein tensor components are
\[
  G_{00}
  = 3\left(\frac{\dot{a}}{a}\right)^2 + 3\frac{k}{a^2},
\qquad
  G_{ij}
  = -\left[
    2\frac{\ddot{a}}{a}
    + \left(\frac{\dot{a}}{a}\right)^2
    + \frac{k}{a^2}
  \right] g_{ij},
\]
with dots denoting time derivatives.

\section{Lamphrodynamic Stress Tensor $\Theta_{\mu\nu}$}
\label{sec:theta-flrw}

From the AKSZ derivation (see Volume~I, \cref{chap:aksz-rsvp}), the lamphrodynamic
stress tensor for homogeneous configurations takes the form
\[
  \Theta_{00}
  = \rho_{\mathrm{lam}}(t),
\qquad
  \Theta_{ij}
  = p_{\mathrm{lam}}(t)\,g_{ij},
\]
where
\begin{equation}
\label{eq:rho-lam}
  \rho_{\mathrm{lam}}
  := \alpha_1 \dot{\SField}^2 + \alpha_2 \SField^2 + \alpha_3\dot{\PhiField}^2,
\end{equation}
\begin{equation}
\label{eq:p-lam}
  p_{\mathrm{lam}}
  := \beta_1 \dot{\SField}^2 + \beta_2 \SField^2 + \beta_3\dot{\PhiField}^2
\end{equation}
for certain coupling constants $\alpha_i,\beta_i$ determined by the AKSZ
action and variational principle. (Explicit formulas may be model‐dependent
and appear in Appendix~I of this volume.)

\begin{remark}
Entropy production $\dot{\SField}>0$ contributes positively to
$\rho_{\mathrm{lam}}$, potentially driving accelerated expansion even in the
absence of a cosmological constant $\Lambda$.
\end{remark}

\section{Modified Friedmann Equations}

Inserting $G_{\mu\nu}$ and $\Theta_{\mu\nu}$ into
\[
  G_{\mu\nu} + \Lambda g_{\mu\nu} + \Theta_{\mu\nu}
  = 8\pi G\,T_{\mu\nu},
\]
with perfect‐fluid matter stress tensor
\[
  T_{00} = \rho, \qquad T_{ij} = p\,g_{ij},
\]
we obtain
\begin{equation}
\label{eq:friedmann1-rsvp}
  3\left(\frac{\dot{a}}{a}\right)^2
  + 3\frac{k}{a^2}
  = 8\pi G\,\rho
  + \Lambda
  + \rho_{\mathrm{lam}},
\end{equation}
\begin{equation}
\label{eq:friedmann2-rsvp}
  -2\frac{\ddot{a}}{a}
  - \left(\frac{\dot{a}}{a}\right)^2
  - \frac{k}{a^2}
  = 8\pi G\,p
  - \Lambda
  + p_{\mathrm{lam}}.
\end{equation}

\section{Effective Equation of State}

Define an effective lamphrodynamic equation‐of‐state parameter
\[
  w_{\mathrm{lam}}
  := \frac{p_{\mathrm{lam}}}{\rho_{\mathrm{lam}}}.
\]
Depending on the signs and magnitudes of $\alpha_i,\beta_i$, one may have
\[
  -1 \le w_{\mathrm{lam}} \le +1,
\]
with $w_{\mathrm{lam}}\approx -1$ corresponding to dark‐energy‐like behavior
and $w_{\mathrm{lam}}\approx 0$ mimicking cold dark matter. In this sense,
lamphrodynamic entropy can \emph{unify} dark energy and dark matter behavior
in a single effective field.

\section{Critical Densities}

Define the critical lamphrodynamic density
\[
  \rho_{\mathrm{crit,lam}}
  := \frac{3}{8\pi G}\left(\frac{\dot{a}}{a}\right)^2
  - \rho - \frac{\Lambda}{8\pi G}.
\]
This characterizes the lamphrodynamic contribution needed to close the
universe ($k=+1$) or accelerate expansion in the flat case ($k=0$). Unlike
in $\Lambda$CDM, critical density may be time‐dependent due to entropy
production.

\section{Discussion}

Lamphrodynamic fields generically modify the cosmic expansion history,
potentially explaining accelerated expansion without a cosmological constant
and altering the interpretation of dark matter. Subsequent chapters explore
these consequences in detail, focusing first on redshift without metric
expansion (\cref{chap:redshift}), then structure formation and dark‐sector
phenomenology.

\clearpage

